\paragraph{Allocation alone does not remove the power cost.}
We test two cheaper alternatives to equal splitting: $\alpha_h=\alpha m_h/m$, and one BH over stratum-specific p-values. For the stated random-role design, conditioning on all stratum counts makes the blocks independent; within-block PRDS then implies pooled PRDS (proof above). Both alternatives control FDR under this design. Neither dominates in power. In the Amazon 80/20 partition with half revelation, attribute-score power is $0.581$ with equal splitting, $0.277$ with size allocation and $0.492$ with pooled stratified ranks. The full random reference gives $0.779$. Changing the allocation alone leaves a substantial power gap in this setting.

\paragraph{When the acquired reference is already biased.}
We next acquire labels with fixed stratum probabilities, using training labels alone to define exposure. At probabilities $(.5,.1)$ for the lower/upper exposure strata, T-Finance graph-score FDP is $0.154$ with pooled biased-reference BH, versus $0.092$ with pooled stratified ranks; powers are $0.818$ and $0.640$. Known-propensity weighted BH has FDP $0.095$ and power $0.785$, but its marginal-rank justification does not supply a joint BH guarantee. Weighted BY supplies a conservative alternative. Figure~\ref{fig:allocation} compares the moderate-acquisition setting on all three graphs; Weibo shows that sparse strata can lose nearly all power. The complete allocation and severity grid is retained in Figure~\ref{fig:allocationfull}.
